\subsubsection{Модель когерентного фотогальванического эффекта}\label{subsubsec:dianov}
\nomenclature{КФЭ}{Когерентный фотогальванический эффект}

В основе модели Дианова и Стародубова лежит представление о том, что фотоиндуцированная решётка квадратичной нелинейности возникает вследствие макроскопического разделения зарядов под действием когерентного фотогальванического эффекта (КФЭ)~\cite{dianov1995photoinduced}. В отличие от ориентационных моделей, использующих оптическое выпрямление, где величина насыщенного постоянного поля не превышает $\sim1$~В/см, КФЭ позволяет достичь напряжённостей электростатического поля $E_{dc}\sim10^4$--$10^5$~В/см, что достаточно для нарушения центросимметричности среды.

При одновременном облучении стекла двумя когерентными пучками с частотами $\omega$ и $2\omega$ возникает интерференция двухфотонного поглощения на частоте накачки и однофотонного поглощения на частоте второй гармоники. В результате появляется когерентный фотогальванический ток, плотность которого феноменологически задаётся в следующем виде~\cite{dianov1995photoinduced}:

\nomenclature[N]{$j_{ph}$~[\si{\ampere\per\meter\squared}]}{Плотность фототока, индуцированного в германосиликатном волокне}
\nomenclature[N]{$\beta$, $\beta_1$, $\ldots$}{Нелинейные коэффициенты, определяющие эффективность соответствующих процессов}

\begin{equation}\label{eq:photocurrent_dianov}
	j_{ph} = \beta E_\omega E_\omega E_{2\omega}^* \exp(\im\Delta k z) + \beta_1 E_\omega E_\omega E_\omega E_{2\omega}^* E_{2\omega}^* \exp(\im\Delta k z) + \ldots + \text{к.с.},
\end{equation}

где $\Delta k = k_{2\omega} - 2k_\omega$~--- фазовая расстройка. Первый член соответствует КФЭ наименьшего (третьего) порядка, последующие — высшим порядкам (пятому, седьмому и т.д.).

Важным вопросом является установление порядка КФЭ при ГВГ.	

Фотогальванический ток приводит к разделению зарядов и накоплению пространственного заряда. Электростатическое поле $E_{dc}$ в среде формируется в результате компенсации фотогальванического тока дрейфовым: $E_{dc} = -j_{ph}/\sigma$, где $\sigma$ — фотопроводимость. Это поле приводит к росту коэффициента квадратичной восприимчивости согласно~\cref{eq:chi2_chi3}, изменение которой, в свою очередь, описывается уравнением


\nomenclature[N]{$A_{\chi}$~[\si{\meter\per\volt\per\second}]}{Скорость индуцирования $\chi^{(2)}$}
\nomenclature[N]{$\tau_d$~[\si{\second}]}{Характерное время релаксации квадратичной нелинейной восприимчивости}
\nomenclature[N]{$\sigma$~[\si{\siemens\per\meter}]}{Удельная электрическая проводимость}

\begin{equation}\label{eq:chi2_growth}
	\dfrac{\partial \chi^{(2)}}{\partial t} = A_{\chi} - \frac{\chi^{(2)}}{\tau_d},
\end{equation}

где $\tau_d = \varepsilon\varepsilon_0/\sigma$; $\varepsilon$ — относительная диэлектрическая проницаемость среды; $\varepsilon_0$ — диэлектрическая проницаемость вакуума; $A_{\chi} \propto \chi^{(3)} j_{ph} / \varepsilon_0$. На начальном участке, когда можно пренебречь релаксацией ($\chi^{(2)}$ мало), рост квадратичной нелинейной восприимчивости линеен, что приводит к квадратичной зависимости сигнала второй гармоники от времени:

\begin{equation}\label{eq:shg_quadratic}
	I_{ph}(2\omega) = \Gamma_d t^2 + \text{const}.
\end{equation}

\nomenclature[N]{$\Gamma_d$~[\si{\watt\per\meter\squared\per\second\squared}]}{Коэффициент квадратичного роста интенсивности второй гармоники во времени}

Параметр скорости $\Gamma_d$ зависит от порядка нелинейного процесса и интенсивностей записывающих волн:

\begin{align}\label{eq:dianov_gamma}
	\Gamma_d \propto &\beta^2 (I_\omega^{\text{write}})^2 I_{2\omega}^{\text{write}} (I_\omega^{\text{read}})^2
	+ \beta_1^2 (I_\omega^{\text{write}})^4 I_{2\omega}^{\text{write}} (I_\omega^{\text{read}})^2+\ldots\\
	\ldots+ &\beta_2^2 (I_\omega^{\text{write}})^2 (I_{2\omega}^{\text{write}})^3 (I_\omega^{\text{read}})^2
	 \beta_3^2 (I_\omega^{\text{write}})^6 I_{2\omega}^{\text{write}} (I_\omega^{\text{read}})^2
	+ \ldots
\end{align}

%В работе \cite{dianov1995photoinduced} для определения порядка КФЭ исследовались зависимости начальной скорости роста сигнала фотоиндуцированной второй гармоники $\Gamma_d$ от интенсивностей записывающего ($I_\omega^{\mathrm{write}}$) и считывающего ($I_\omega^{\mathrm{read}}$) излучений в двух режимах: при совмещённых и при разнесённых во времени импульсах. В случае несовмещённых импульсов показатели степеней составили $6$ по $I_\omega^{\mathrm{write}}$ и $2$ по $I_{\omega}^{\mathrm{read}}$, что соответствует участию шести когерентных фотонов в процессе. При совмещении импульсов за счёт заселения промежуточных состояний считывающим излучением наблюдалось резкое (примерно в 100 раз) увеличение скорости роста и изменение показателей степени до $2$ по $I_\omega^{\mathrm{write}}$ и $6$ по $I_{\omega}^{\mathrm{read}}$. Эти зависимости соответствуют наименьшему порядку КФЭ, в котором участвуют три когерентных фотона ($\omega + \omega + 2\omega$), а вклад считывающего излучения сводится к двухфотонному заселению промежуточного состояния. Согласно модели авторов, в таком режиме начальная скорость роста сигнала ГВГ определяется выражением

В работе \cite{dianov1995photoinduced} для определения порядка КФЭ исследовались зависимости начальной скорости роста сигнала фотоиндуцированной второй гармоники $\Gamma_d$ от интенсивностей записывающего ($I_\omega^{\mathrm{write}}$) и считывающего ($I_\omega^{\mathrm{read}}$) излучений в двух режимах: при совмещённых и при разнесённых во времени импульсах. Эти зависимости соответствуют наименьшему порядку КФЭ, в котором участвуют три когерентных фотона ($\omega + \omega + 2\omega$), а вклад считывающего излучения сводится к двухфотонному заселению промежуточного состояния. Согласно модели авторов, в таком режиме начальная скорость роста сигнала ГВГ определяется выражением

\begin{equation}\label{eq:gamma_dianov}
	\Gamma_d \propto \beta^2 (I_\omega^{\mathrm{write}})^2 I_{2\omega}^{\mathrm{write}} (I_\omega^{\mathrm{read}})^2,
\end{equation}

 Наличие множителя $(I_\omega^{\mathrm{read}})^2$ отражает двухфотоннное заселение промежуточного состояния считывающим излучением, а произведение $(I_\omega^{\mathrm{write}})^2 I_{2\omega}^{\mathrm{write}}$ соответствует когерентной интерференции двухфотонной ионизации накачкой и однофотонной ионизации второй гармоникой.
 
 
\todo{Критика Хмели!! Непонятно какие величины. 100\% эффективность}

в выводе тоже критика


\todo{Таким образом, модель Дианова и Стародубова объясняет формирование решётки $\chi^{(2)}$ макроскопическим разделением зарядов под действием когерентного фотогальванического тока. Она успешно описывает степенные зависимости скорости роста от интенсивностей, ускорение записи при наличии затравочной второй гармоники и эффекты оптического стирания решётки. Однако, как будет показано в следующем разделе, существует альтернативный подход — модель экситонов с переносом заряда Антонюка, которая предлагает иной микроскопический механизм возникновения поля $E_{dc}$ и предсказывает ряд эффектов (в частности, самовыключение ГВГ), не рассматриваемых в рамках фотогальванической модели.}